\subsection{Konteks Aplikasi}

PasarKita adalah aplikasi marketplace digital untuk produk UMKM. Aplikasi ini mempertemukan tiga aktor utama: pembeli, penjual, dan admin. Pembeli menggunakan aplikasi untuk melihat katalog, mencari produk, memasukkan produk ke keranjang, melakukan checkout, dan melihat status pesanan. Penjual menggunakan aplikasi untuk mengelola produk, stok, pesanan, serta dashboard toko. Admin menggunakan aplikasi untuk memantau user, produk, order, analytics, dan kondisi integrasi.

Marketplace seperti PasarKita membutuhkan algoritma karena data yang ditampilkan tidak cukup hanya disimpan dan diambil apa adanya. Produk perlu direkomendasikan, diurutkan, dicari, dan dihitung biayanya. Recommender system membantu pengguna memilih item ketika pilihan terlalu banyak \parencite{Resnick1997}. Dari sisi organisasi, penggunaan teknologi informasi juga relevan untuk membantu proses adopsi digital pada level usaha atau organisasi \parencite{Oliveira2011}. Karena itu, PasarKita cocok dijadikan objek analisis penerapan algoritma dasar.

\subsection{Fokus Analisis}

Laporan ini hanya membahas fitur yang sudah memiliki penerapan algoritma jelas di dokumentasi dan kode aplikasi. Fokusnya adalah sebagai berikut:

\begin{itemize}[leftmargin=*]
    \item Algoritma Greedy pada rekomendasi produk, sorting produk, dan fee marketplace.
    \item String Matching/KMP pada pencarian produk buyer dan seller.
    \item String Matching multi-field pada pencarian pesanan.
    \item Potongan kode, pseudocode, dan contoh output dari penerapan algoritma.
\end{itemize}

Laporan ini tidak mengklaim bahwa PasarKita sudah memakai machine learning recommendation. Referensi Epsilon Greedy digunakan sebagai catatan pengembangan, bukan sebagai implementasi saat ini \parencite{Umami2021}.

\subsection{Fitur yang Dianalisis}

Fitur yang dianalisis dipilih karena langsung berhubungan dengan kebutuhan marketplace.

\begin{table}[htbp]
    \caption{Fitur PasarKita yang Dianalisis}
    \label{tab:fitur-dianalisis}
    \small
    \begin{tabularx}{\textwidth}{@{} >{\raggedright\arraybackslash}X >{\raggedright\arraybackslash}X >{\raggedright\arraybackslash}X @{}}
    \toprule
    \textbf{Fitur} & \textbf{Masalah yang Diselesaikan} & \textbf{Algoritma} \\
    \midrule
    Rekomendasi produk & Pembeli butuh produk yang relevan dari banyak pilihan & Greedy scoring \\
    Sorting produk & Produk perlu ditampilkan berdasarkan performa penjualan atau rating & Greedy ranking \\
    Fee marketplace & Total checkout perlu dihitung cepat dan konsisten & Greedy satu langkah \\
    Pencarian produk & Buyer dan seller perlu menemukan produk berdasarkan keyword & KMP/String Matching \\
    Pencarian pesanan & Seller perlu mencari order berdasarkan beberapa field & String Matching multi-field \\
    \bottomrule
    \end{tabularx}
\end{table}

\subsection{Sumber Data Analisis}

Analisis dilakukan dari tiga sumber utama. Pertama, dokumentasi algoritma pada folder \path{docs/algorithms}. Kedua, source code aplikasi yang berkaitan dengan rekomendasi, fee, pencarian produk, dan pencarian pesanan. File utama yang ditinjau adalah modul rekomendasi frontend, utilitas fee, utilitas KMP, service produk, dan service order. Ketiga, referensi yang sudah dikumpulkan dalam file BibTeX \path{src/My Collection.bib}.
